% Philippica I (philippic-1) --- English translation
% Translated by Claude (Anthropic) from the Latin (Perseus).
% Style guide: STYLE.md.

\ciceroSection{1} Before I speak, senators, of what I think must be said in the present state of the commonwealth, I shall set out for you briefly the reasons both for my going away and for my coming back. As long as I had any hope that the commonwealth had at last been recalled to your counsel and authority, I held that I must stay in place, as if on a consular and senatorial watch. Nor indeed did I ever depart anywhere, nor did I take my eyes off the commonwealth, from the day on which we were summoned together to the temple of Tellus. In that temple, so far as it lay in my power, I laid the foundations of peace, and renewed the old precedent of the Athenians; I even pressed into service the Greek word that city had once used in calming its civil discord \textit{[Greek: amnestia]}, and I urged that all memory of our dissensions ought to be wiped out by an everlasting forgetfulness.

\ciceroSection{2} Brilliant then was the speech of Marcus Antonius, and his goodwill outstanding too; peace, in the end, was confirmed through him and through his children with the most distinguished of our citizens. And the rest answered to such beginnings. To those deliberations on the commonwealth which he held at his own house he invited the leading men of the state; to this order he brought the best of measures; nothing was then found in the records of Gaius Caesar except what was already known to all; he answered with complete consistency the questions that were put to him.

\ciceroSection{3} Had any exiles been brought back? One, he said, and besides him no one. Had any exemptions been granted? ``None,'' he answered. He even wanted us to agree to the proposal of Servius Sulpicius, a most distinguished man, that no tablet should be set up after the Ides of March bearing any decree or grant of favour of Caesar's. I pass over many things, and brilliant ones; for my speech hurries on to one singular act of Marcus Antonius. The dictatorship, which had by then usurped the force of regal power, he rooted out of the commonwealth entirely; on that we did not even cast our votes. He brought in a senatorial decree, drawn up as he wished it to be passed; when it had been read out, we followed his authority with the highest enthusiasm and gave him thanks, by a decree of the Senate, in the amplest terms.

\ciceroSection{4} A kind of light seemed to have been let in upon us, now that not only the kingship that we had endured but even the fear of kingship had been lifted away, and a great pledge had been given by him to the commonwealth that he wished it to be a free state, when he had rooted out of the commonwealth entirely the very name of dictator --- a name that had often been legitimate --- because the perpetual dictatorship was still fresh in memory.

\ciceroSection{5} The Senate, a few days later, was freed from the danger of slaughter; the hook was driven into that runaway slave who had usurped the name of Marius. And all of this in common with his colleague; the rest, again, belongs properly to Dolabella --- and would, I believe, have been shared in common, had his colleague not been away. For while a measureless evil was creeping through the city and spreading wider day by day, while the very men who had given that unburied corpse its burial were making a funeral pyre of it in the forum, and while every day, more and more, abandoned wretches with their like among the slaves were threatening the houses and temples of the city, such was Dolabella's chastisement, both of the audacious and criminal slaves and of the foul and impious freedmen, and such his overthrow of that accursed column, that I find it astonishing how widely the time that followed has differed from that one day.

\ciceroSection{6} For lo, on the Kalends of June, on which they had given notice for us to be present, everything was changed: nothing through the Senate, much business of weight through the people, and that with the people absent and against its will. The consuls-elect kept saying that they did not dare come into the Senate; the liberators of their country were kept away from the city whose servile yoke they had thrown off her neck --- though the consuls themselves in their public speeches and in all their conversation kept praising them. The so-called veterans, for whom this order had taken the most diligent precautions, were being stirred up not to the protection of what they had but to the hope of new plunder. Since I preferred to hear of these things rather than to see them, and since I had the right of a free legation, I went away in the resolve to be present on the Kalends of January, which seemed likely to be the beginning of the Senate's being convened.

\ciceroSection{7} I have set out, senators, the reasoning of my departure: I shall now briefly set out that of my return, which has more in it to surprise. When I had avoided Brundisium and that worn road to Greece, and not without cause, on the Kalends of August I came to Syracuse, since the crossing into Greece from that city was well spoken of: yet that city, most closely tied to me as it is, for all its wishing to keep me longer than a single night, could not. I was afraid that my sudden arrival to my friends there might bring some shadow of suspicion upon them, if I lingered. When the winds had carried me from Sicily as far as Leucopetra, which is a promontory of the Regian territory, from that place I embarked to make the crossing; but not very far out, I was driven back by a southerly gale to the very spot from which I had embarked.

\ciceroSection{8} When night had fallen and I had stayed at the country house of Publius Valerius, my companion and intimate, and the next day, still under the same roof, was waiting for a wind, many of the townsmen of Regium came to see me, some of them fresh from Rome: from these I first received Marcus Antonius's address to the assembly --- which so pleased me, on reading it, that I began for the first time to think of turning back. Not long after, the edict of Brutus and Cassius was brought to me, which to me, perhaps because I love them even more for the commonwealth's sake than for friendship's, seemed full of equity. They added besides --- for it commonly happens that men who wish to bring good news add something of their own to make what they announce more cheerful --- that there would be an accommodation: the Senate would be in full session on the Kalends; Antony, having sent his bad advisers packing and given up the Gallic provinces, would come back to the Senate's authority.

\ciceroSection{9} Then in truth I was set on fire with such eagerness to return that no oars and no winds could satisfy me, not that I thought I should fail to reach Rome in time, but that I might not give thanks to the commonwealth more slowly than I longed to. Sailing rapidly down to Velia, I saw Brutus: with what grief on my part, I do not say. It seemed shameful to me to dare to come back into the city out of which Brutus was withdrawing, and to wish to be safe in a place where he could not be. Yet I did not find him stirred as I myself was. Rather, upraised by his consciousness of his very great and most beautiful deed, he was complaining of nothing for himself, much for you.

\ciceroSection{10} And from him I first learned what speech Lucius Piso had made in the Senate on the Kalends of August: a speech which, although he was too little --- this very thing I had heard from Brutus --- supported by those whose duty it was, yet by Brutus's witness --- and what can be of greater weight than that? --- and by the testimony of everyone I have seen since, seemed to me to have won great glory. To follow him, then, I made haste --- him whom those present had not followed --- not that I expected to accomplish anything (for I neither hoped for it nor could promise it), but so that, if anything in the way of human chance should befall me (and many things appear to be hanging over us beyond nature and beyond fate), I might still leave the voice of this day as a witness to the commonwealth of my unbroken goodwill towards her.

\ciceroSection{11} Since I am confident, senators, that the reasoning of both decisions has been approved by you, before I begin to speak of the commonwealth I shall briefly complain of yesterday's affront from Marcus Antonius --- of whose friend I am, and which, for some service of his, I have always made open profession that I ought to be. What, then, was the cause that yesterday I should be so harshly compelled into the Senate? Was I alone absent? or have you not often been less than full in number? or was a matter being transacted that called for even the sick to be brought in? Hannibal, I suppose, was at the gates, or the question of peace with Pyrrhus was being debated --- the very business to which, the record says, that famous Appius was brought in, blind and aged.

\ciceroSection{12} The motion was concerning supplications --- a kind of business in which senators do not usually fail to attend. For they are compelled not by sureties but by their regard for the men whose honour is in question; just as happens when a triumph is brought forward. The consuls accordingly take no trouble about it, so that it is almost a senator's free choice whether or not to be there. Since this practice was known to me, and since I was weary from the road and unhappy with myself, I sent, in the name of our friendship, a man to tell him as much. But he, with you to hear him, said that he would come to my house with builders. This was both too angry and very intemperate. For what is the offence of which this is so great a penalty, that he should dare to say in this order that he would use the public workmen to pull down a house publicly built by decree of the Senate? Who has ever compelled a senator's attendance with so great a fine? Or what penalty is there beyond surety or fine? But had he known the opinion I was going to deliver, he would surely have remitted something of the harshness of his summoning.

\ciceroSection{13} Or do you suppose, senators, that I should have voted, as you reluctantly voted, that funeral rites should be mingled with thanksgivings, that inexpiable religious offences should be brought into the commonwealth, that thanksgivings should be decreed for a dead man? I will not say for whom. Let it have been that famous Lucius Brutus, who both himself liberated the commonwealth from regal mastery and propagated his stock to a like virtue and a like deed almost into the five-hundredth year: I could not be brought to join any dead man with the worship of the immortal gods --- so that public supplication should be made to one whose tomb stands somewhere with rites of the dead performed at it. For my part I should have given an opinion such that I could easily defend it before the Roman people if some heavier disaster overtook the commonwealth --- if there were war, if pestilence, if famine; some of which are already upon us, others of which I fear are hanging over us. But may the immortal gods pardon this, I pray, both to the Roman people, which does not approve it, and to this order, which decreed it against its will.

\ciceroSection{14} But what? On the other evils of the commonwealth, is one allowed to speak? Yes, I am allowed, and I shall always be allowed, to uphold the dignity of my office and to think nothing of death. Only let me have the power to come into this place: I shall not refuse the peril of speaking. And I wish, senators, that I had been able to be present on the Kalends of August! Not that anything could have been accomplished, but only that there might not be found, as on that day there was, one consular alone worthy of that office and worthy of the commonwealth. From which fact I take a great grief --- that men who had enjoyed the most ample favours of the Roman people did not follow Lucius Piso as the leader of the best of opinions. Did the Roman people make us consuls for this --- that, set on the highest step of dignity, we should count the commonwealth as nothing? Not one man assented to the consular Lucius Piso, not by voice, not even by a look.

\ciceroSection{15} What, the devil, is this voluntary slavery? Some slavery there was that could not be helped; and I do not ask this kind of thing of all those who deliver their opinion in the place of consulars. The case is different of those whose silence I forgive; different again of those whose voice I miss. With grief I see them coming under the suspicion of the Roman people --- not of fear (which would itself be shameful) but that each, for one cause or another, is failing his own dignity. For which reason I both give and feel the greatest thanks to Piso, who did not think what he could accomplish in the commonwealth but what he himself ought to do. And next I ask of you, senators, that even if you will be less bold in following my speech and my authority, you will still hear me kindly, as you have done up to now.

\ciceroSection{16} First, then, I hold that Caesar's acts must be upheld --- not that I approve them (for who can do that?) but because I judge the highest priority should be given to peace and quiet. I could wish that Marcus Antonius were present, only without his retinue of advocates --- but, as I take it, he is allowed to be unwell, where yesterday on his account that was not allowed to me --- and were here to teach me, or rather to teach you, senators, in what manner he himself would defend Caesar's acts. Are they, set forth in those little notebooks and handwritten papers and memoranda, produced under his sole authority --- not even produced, but merely talked of --- to be the firm acts of Caesar; and the things which Caesar himself cut into bronze, in which he willed that the orders of the people and the laws should be made perpetual, to be counted for nothing?

\ciceroSection{17} For my part I judge that nothing is so much in Caesar's acts as Caesar's laws. Is it that, if he promised something to any man, that shall stand fixed, the very thing he could not do? --- as he failed to make good many promises to many men: of which still many more have been discovered since his death than were ever conferred and given as favours through all the years of his life. But I do not change these or move them; I defend his brilliant acts with the highest enthusiasm. Would that the money in the temple of Ops were untouched! Bloody money, indeed, but in these times, since it is not being given back to those whose it is, money the commonwealth needs. Though let even that be poured out, if it stood so in his acts.

\ciceroSection{18} Is there anything that can so properly be called an act of one who has held office, with power and command, in a civilian capacity in the commonwealth, as a law? Ask for the acts of Gracchus: the Sempronian laws will be produced; ask for those of Sulla: the Cornelian laws. What of Pompey's third consulship? On what acts did it stand? On laws, surely. If you asked Caesar himself what he had done in the city and in the toga, he would answer that he had carried many laws, and brilliant ones, but the handwritten papers --- he would either change them, or not produce them, or, if he had produced them, would not number those things among his acts. But on these very things I yield; in some matters I even shut my eyes; but in the greatest matters, that is, in the laws, that Caesar's acts should be undone I cannot endure.

\ciceroSection{19} What law was better, more useful, more often demanded even when the commonwealth was at its best, than that praetorian provinces should not be held for more than one year, nor consular provinces for more than two? Does it seem to you that, with this law overthrown, Caesar's acts can be upheld? Again, in the law that has been promulgated concerning a third decury (\textit{tertia decuria} --- a new panel of jurors), are not all of Caesar's juror-laws (\textit{leges iudiciariae}) being undone? And do you defend Caesar's acts who overturn his laws? Unless, perhaps, if he set down some memorandum in a notebook, that is to be reckoned among his acts and, however unjust and useless it may be, defended; and what he carried before the people in the centuriate assembly is not to be held among the acts of Caesar.

\ciceroSection{20} But what is this third decury? ``Of centurions,'' he says. What? Was that order not open to a place on the jury under the Julian law (\textit{lex Iulia}), and earlier under the Pompeian (\textit{lex Pompeia}) and the Aurelian (\textit{lex Aurelia})? ``A property qualification was prescribed,'' he says. Not for the centurion alone but for the Roman knight as well; and so the bravest and most honourable of men who have led companies of the line both sit on juries and have done so. ``I am not asking after those,'' he says. ``Whoever has led a company, let him be a juror.'' Yet if you proposed that anyone at all who had served in the cavalry --- which is a more polished thing --- should be a juror, you would persuade no one; for in a juror both fortune and dignity ought to be looked at. ``I am not asking after those things,'' he says. ``I add even jurors taken from the rank and file out of the Legion of the Larks. For otherwise our men say they cannot be safe.'' What an insulting honour for those you summon to judge against their expectation! For this is what the law is in fact pointing to --- that those should judge in the third decury who would not dare to judge in freedom. And in this what an error --- by the immortal gods! --- in those who thought up such a law! For just as each man shall seem most disreputable, just so most gladly shall he wash off his disreputability by the strictness of his judging, and shall labour to seem worthy rather of the respectable decuries than rightly cast into the contemptible.

\ciceroSection{21} Another law has been promulgated --- that those condemned of violence and of treason should be allowed to appeal to the people, if they wish. Is this, then, a law, or rather the dissolving of all laws? For who is there today whose interest it is that such a law should stand? No one is under indictment by those laws, no one whom we suppose ever will be. For things done by arms will surely never be called into a court of law. ``But it is a popular measure.'' Would that you had wished for anything to be popular! For all citizens already agree, with one mind and one voice, about the safety of the commonwealth. What then is this desire to carry a law that has the highest disgrace in it and no favour? For what is more disgraceful than that the man who has by violence lessened the majesty of the Roman people, when condemned by judgment, should return to that very violence on whose account he was rightly condemned?

\ciceroSection{22} But why do I argue any further about the law? As if the point were that anyone should appeal: the point, the thing being carried, is that no one shall ever in any way be put on trial under those laws. For what prosecutor will be found mad enough to be willing, with the defendant condemned, to expose himself to a hired mob? or what juror to dare condemn a defendant, only to be himself dragged off at once before the hired claqueurs? It is not, then, the right of appeal that this law grants; rather, two of the most salutary laws and the standing courts they create are being abolished. What else is this than to incite our young men to wish to be turbulent, seditious, ruinous citizens? To what bane will the rage of a tribune not be driven, with these two inquiries into violence and treason taken away?

\ciceroSection{23} What further? Caesar's laws are being repealed --- laws that command that the man condemned of violence, and likewise the man condemned of treason, shall be barred from fire and water. When the right of appeal is given to such men, are not Caesar's acts cut down? Yet I, senators, who never approved those laws, judged still that they should be kept up for the sake of concord, in such a way that I thought not only laws Caesar carried while living must not be undermined at this time, but not even those which you see produced and posted after his death.

\ciceroSection{24} Men brought back from exile by a dead man; citizenship granted not only to individuals but to entire peoples and provinces by a dead man; tax-revenues abolished by exemptions without limit by a dead man. So we are defending these things, brought from his house on the single but supreme authority of one man; and the laws which he himself, with us looking on, read out, proclaimed, and carried --- in whose carrying he gloried, and in whose laws he held the commonwealth to be held together --- the laws on provincial commands, on the courts, those laws, I say, the laws of Caesar, we who defend Caesar's acts are to think must be overthrown?

\ciceroSection{25} And yet of these laws that have been promulgated we can at least complain: of those which are said already to have been carried, even that has not been allowed. For they were carried without any promulgation at all, before they were drafted. But I ask what reason there is why I, or any of you, senators, should fear bad laws under good tribunes of the plebs. We have men ready to interpose their veto; ready to defend the commonwealth on grounds of religion: we ought to be free of fear. ``What vetoes do you bring up to me,'' he says, ``what religious grounds?'' Those, of course, on which the safety of the commonwealth rests. ``We have no regard for such things; we count them too old and silly: the forum shall be fenced off; every approach shall be closed; armed men shall be posted in many places.'' What then?

\ciceroSection{26} Whatever is done in that fashion --- shall that be law? And you will give orders, I take it, to cut into bronze those formulaic words: ``The consuls duly asked the people'' --- is this the right of asking that we inherit from our ancestors? ``and the people duly voted.'' What people? the one that was shut out? By what right? By the right that everything has been swept away by force and arms? I am speaking of what is to come, since it is the part of friends to speak in advance of what can be avoided: if these things shall not be done, my speech will stand refuted. I am speaking of laws promulgated, on which you still have a free hand; I am pointing out the faults: take them away! I am giving notice of violence: take away the arms!

\ciceroSection{27} Nor will it become you, Dolabella, to be angry with me when I speak for the commonwealth. Though I do not think that you, at any rate, will do so --- I know your easy temper --- but they say your colleague has become an angry man in this fortune of his which seems good to him --- though it would seem to me, not to speak of him more harshly, more fortunate if he imitated his grandfather's and his uncle's consulships. I see, however, what an odious thing it is to have the same man angry and armed, especially when there is such impunity for swords; but I shall propose terms which I think fair, and which I do not suppose Marcus Antonius will refuse. If I shall have said anything against his life or his manners by way of insult, I shall not deprecate that he be my bitterest enemy; but if I keep the practice I have always had in public matters, that is, if I speak freely what I think about the commonwealth, in the first place I beg him not to be angry; and next, if I do not obtain this, I ask that he be angry as a citizen. Let him have his arms, if it must be, as he says, for the sake of his own defence: only let those arms not harm men who shall have said for the commonwealth what seemed right to them. What can be said fairer than this demand?

\ciceroSection{28} But if --- as has been said to me by some of his intimates --- any speech delivered against his will offends him deeply, even if it contains no insult, we shall bear with the nature of a friend. But those same men also say this to me: ``You, an opponent of Caesar's, will not be allowed what was allowed to Piso his father-in-law''; and at the same time they warn me of something I shall be on guard against: ``Sickness will be no more lawful an excuse for not coming into the Senate than death.''

\ciceroSection{29} But, by the immortal gods --- for as I look at you, Dolabella, who are most dear to me, I cannot keep silent about the error of you both --- I do believe that you nobles, looking to something great, have aspired not to money (which some, who are too credulous, suspect, and which the noblest and most illustrious men have always despised), nor to violent power, nor to a domination too heavy for the Roman people to bear, but to the affection of the citizens and to glory. Now glory is the praise of right deeds and the great repute of services rendered to the commonwealth, which is approved by the witness of the best men no less than by that of the multitude.

\ciceroSection{30} I would have said, Dolabella, what the fruit of right deeds is, if I did not see that you, beyond others, have for a little while had experience of it. What day in your life can you remember as having dawned more joyfully than that when, the forum purged, the throng of the impious scattered, the ringleaders of the crime punished, the city freed from fire and the fear of slaughter, you returned to your house? Of what order, of what class, of what fortune were the zealous demonstrations of praise and congratulation which did not offer themselves to you? Indeed, even to me did the good men give thanks --- for by my advice they thought you were guiding yourself in these affairs --- and they congratulated me in your name. Remember, I beg, Dolabella, that consensus of the theatre, when all, forgetting the things on whose account they had been offended with you, made it plain that by your recent kindness they had cast off the memory of their old grievance.

\ciceroSection{31} This, Publius Dolabella --- I speak in great grief --- this, I say, were you able to lay aside in calm of mind, so great a dignity? And you, Marcus Antonius --- for I address you in your absence --- do you not set that one day on which the Senate was in the temple of Tellus above all these months in which certain men, very much my opposites in opinion, suppose you to be happy? What a speech yours was that day, on concord! From what fear was the Senate, from what trouble was the state freed by you that day, when, your enmities laid aside, forgetful of the auspices announced from yourself as augur of the Roman people, you willed that on that first day your colleague should be your colleague --- when your little son was sent up by you into the Capitol as a pledge of peace!

\ciceroSection{32} On what day was the Senate more joyful, on what the Roman people --- which indeed in no assembly was ever more thickly thronged? At last we seemed to have been freed by the bravest of men, because, as they had willed, peace was following liberty. On the next day, the second, the third, and finally on all the days that followed, you did not cease to bring each day, as it were, some gift to the commonwealth; but the greatest was that you abolished the name of dictator. This brand was burnt by you --- by you, I say --- upon dead Caesar to his perpetual disgrace. For just as on account of the single crime of Marcus Manlius, by a decree of the Manlian clan, no patrician Manlius is permitted to be called Marcus, so you, on account of the hated name of one dictator, rooted out the name of dictator from the commonwealth entirely.

\ciceroSection{33} Could it be, when you had done such great things for the safety of the commonwealth, that you were dissatisfied with your fortune, your eminence, your renown, your glory? Whence, then, on a sudden, so great a change? I cannot be brought to suspect that you were caught by money. Let each man say what he likes; one is not bound to believe. For I have never known anything sordid in you, anything low. Though, to be sure, intimates of one's own household do sometimes deprave a man; but I know your firmness. And would that you could have avoided, as you avoided the fault, even the suspicion! What I fear more is this --- that, in ignorance of the true road of glory, you may think it glorious to be more powerful, you alone, than all men together, and you may prefer to be feared by your countrymen than loved. If that is your thought, you are wholly ignorant of the road to glory. To be cherished as a fellow citizen, to deserve well of the commonwealth, to be praised, attended on, loved --- that is glorious; but to be feared and an object of hatred is hateful, detestable, weak, and falling.

\ciceroSection{34} As we see, even in the play, was ruinous to that very man who said \textit{``oderint, dum metuant''} --- \textit{let them hate, so long as they fear}. Would, Marcus Antonius, that you might remember your grandfather! And indeed you have heard much about him from me, and very often. Do you suppose that he wished to deserve immortal renown by being feared on account of his licence in bearing arms? That was a life, that was a happy fortune --- to be the equal of others in liberty, the first in dignity. So that, to pass over your grandfather's prosperous years, I should choose his most bitter last day rather than the despotism of Lucius Cinna, by whom he was most cruelly killed.

\ciceroSection{35} But why should I sway you with words? For if the end of Gaius Caesar cannot bring it about that you would rather be loved than feared, nothing anyone says will profit or prevail. Those who think him to have been happy are themselves wretched. No man is happy who lives on such terms that he can be killed not only with impunity but even with the greatest glory to his killer. Therefore, bend yourself, I beg, and look back to your ancestors, and so steer the commonwealth that your countrymen shall rejoice that you were born: without which no man can be either happy or cherished or pleasing at all.

\ciceroSection{36} You both have many judgments from the Roman people, by which I bear it very ill that you are not sufficiently moved. What of the shouts of countless citizens at the gladiatorial games? What of the people's verses? What of the boundless applause at Pompey's statue? What of the two tribunes of the plebs who stand against you? Do these things not signify, beyond believing, the agreement of the whole Roman people? What? Did the applause at the Apollinarian games, or rather the testimonies and judgments of the Roman people, seem too small to you? O happy were those men who, when they themselves were not allowed to be present because of the force of arms, were yet present, and clung to the marrow and inward parts of the Roman people! Unless, perchance, you supposed that it was Accius who was being applauded then, and that the palm was being given in his sixtieth year and not to Brutus --- Brutus, who in this way was absent from his own games, so that at that most lavishly appointed spectacle the Roman people might offer their zeal to a man not present, and might soothe their longing for their liberator with unbroken applause and shouting.

\ciceroSection{37} For my part I am one who has always despised that kind of applause when it was being offered to demagogue politicians; and yet when this same applause comes from the highest and middling and lowest, in short from all together, and when the men who used to follow the people's agreement flee from it, I count it not applause but a judgment. But if these things, which are the gravest, seem light to you, do you despise this too --- that you have felt how dear the life of Aulus Hirtius has been to the Roman people? For it was enough that he was approved by the Roman people, as he is; pleasing to his friends, in which he surpasses all; dear to his own people, to whom he himself is dearest: yet for whom do we remember so great a solicitude of the good men, so great a fear in all? Surely for no one.

\ciceroSection{38} What then? By the immortal gods, do you not interpret what this means? What? Do you not suppose that those, in whose eyes the life of men they hope will look after the commonwealth is so dear, are thinking about your life? I have gathered, senators, the fruit of my return, since I have both said things such that, whatever the outcome might be, there should stand a witness of my constancy, and I have been kindly and attentively heard by you. If this power shall be given to me more often without peril to me or to you, I shall use it: if not, so far as I can, I shall keep myself in reserve, not so much for my own sake as for the commonwealth's. For me what I have lived is roughly enough, whether for years or for glory: if anything is added to it, it shall be added not so much to me as to you and to the commonwealth.
